\chapter*{Conclusion générale}
\addcontentsline{toc}{chapter}{Conclusion générale}

L'objectif de ce mémoire était d'explorer les liens entre deux notions a priori distinctes : les surfaces elliptiques, issues de la géométrie algébrique complexe, et les fibrations de Lefschetz, relevant de la topologie différentielle en dimension 4. Le fil conducteur de cette étude a été la notion de monodromie, qui permet de traduire des questions géométriques en problèmes combinatoires portant sur des factorisations dans le groupe modulaire $SL(2,\Z)$.

Nous avons commencé par rappeler les fondements de la théorie des surfaces elliptiques, en insistant sur la classification de Kodaira des fibres singulières. Cette classification, qui constitue l'un des aboutissements les plus remarquables de la géométrie des surfaces complexes, fournit un inventaire complet des configurations possibles de composantes irréductibles d'une fibre singulière, chacune étant caractérisée par son graphe dual et ses contributions aux invariants globaux de la surface. Les exemples des surfaces $E(1)$, des surfaces K3 et des surfaces de Dolgachev ont illustré la diversité des comportements possibles.

Nous avons ensuite introduit les fibrations de Lefschetz, en montrant comment elles émergent naturellement des surfaces elliptiques par perturbation des fibres singulières. La définition du groupe de mapping class et des twists de Dehn a permis de formaliser la notion de monodromie, qui associe à tout lacet dans la base un difféomorphisme de la fibre régulière. Dans le cas du genre 1, l'identification du groupe de mapping class avec $SL(2,\Z)$ a établi un pont fondamental entre la topologie des surfaces et l'arithmétique.

L'étude systématique de la monodromie comme outil de classification a fait l'objet du troisième chapitre. Nous y avons montré qu'une fibration de Lefschetz sur la sphère est entièrement déterminée par une factorisation de l'identité en produit de twists de Dehn, et que deux factorisations reliées par des mouvements de Hurwitz correspondent à la même fibration. Dans le cas du genre 1, le problème se ramène à l'étude des factorisations de l'identité dans $SL(2,\Z)$ en conjugués de la matrice unipotente $U$. Les travaux de Moishezon, Kas et Matsumoto ont établi une classification complète pour ce cas, tandis que la formule de Meyer a fourni un lien direct entre la monodromie et la signature.

Le dernier chapitre a été consacré à l'étude détaillée d'un exemple concret, la surface $E(1)$. Construite à partir d'un pinceau de cubiques dans le plan projectif éclaté en ses neuf points de base, cette surface elliptique possède douze fibres singulières de type $I_1$. La fibration de Lefschetz associée a une monodromie donnée par douze twists de Dehn, dont le produit est l'identité. Par des mouvements de Hurwitz, cette factorisation se ramène à la relation $(UV)^6 = I$, qui n'est autre que la relation fondamentale du groupe modulaire $PSL(2,\Z)$. Cet exemple a également servi de point de départ pour présenter les généralisations que constituent les surfaces $E(n)$ et les surfaces de Dolgachev.

Au terme de ce parcours, plusieurs enseignements peuvent être tirés. D'une part, l'approche par la monodromie s'avère extrêmement efficace pour réduire des questions géométriques complexes à des problèmes combinatoires maniables. Elle permet de calculer des invariants topologiques fins comme la signature, et fournit un langage unifié pour aborder des objets issus de traditions mathématiques différentes. D'autre part, cette approche a ses limites : la structure complexe, les périodes, le $j$-invariant sont perdus dans le passage à la monodromie topologique, et toutes les factorisations ne correspondent pas nécessairement à des surfaces elliptiques complexes.

Les perspectives ouvertes par ce travail sont nombreuses. Du côté des fibrations de genre supérieur ($g \ge 2$), la classification reste largement incomplète, la complexité du groupe de mapping class rendant le problème autrement plus difficile que pour le genre 1. Du côté de la géométrie symplectique, le théorème de Donaldson assure que toute variété symplectique de dimension 4 admet une fibration de Lefschetz, ce qui fait de ces objets un outil central pour l'étude des structures symplectiques. Du côté de la théorie de Hodge, le lien entre monodromie topologique et variations de structures de Hodge reste à explorer plus avant. Enfin, les questions liées à l'existence de sections, à la détection des fibres multiples ou à l'interprétation des invariants de Seiberg-Witten en termes de monodromie constituent autant de directions de recherche actuelles.

Ce mémoire n'a pu qu'effleurer la richesse d'un sujet qui, depuis les travaux fondateurs de Kodaira et Lefschetz jusqu'aux développements contemporains liés à la symétrie miroir et aux invariants de Donaldson, n'a cessé de se renouveler. Il aura du moins, nous l'espérons, donné au lecteur les clés pour aborder la littérature spécialisée et, qui sait, pour contribuer à son tour à l'édifice.